A short companion to the main results that makes explicit every
statistical choice a reviewer might question. Each subsection answers
one specific question that came up during the work and is therefore
worth anticipating. The numeric justifications reference the same
per-iteration CSVs archived under
\href{../experiments/}{\texttt{../experiments/}} and summarised in
\url{evaluation.md}.

\begin{center}\rule{0.5\linewidth}{0.5pt}\end{center}

\section{1. Why n = 20 (E14/E15) and n = 40
(E16/E18)?}\label{why-n-20-e14e15-and-n-40-e16e18}

Let \texttt{σ} be the sample standard deviation of per-iteration times
and \texttt{SE\ =\ σ/√n} the standard error of the mean. To support a
claim that two means differ by Δ with 95 \% confidence, we want
\texttt{\textbar{}Δ\textbar{}\ ≫\ 2\ ×\ SE}.

Observed \texttt{σ} values for the dominant measurements, from the
archived CSVs:

{\def\LTcaptype{none} % do not increment counter
\begin{longtable}[]{@{}
  >{\raggedright\arraybackslash}p{(\linewidth - 6\tabcolsep) * \real{0.2000}}
  >{\raggedleft\arraybackslash}p{(\linewidth - 6\tabcolsep) * \real{0.2667}}
  >{\raggedleft\arraybackslash}p{(\linewidth - 6\tabcolsep) * \real{0.2667}}
  >{\raggedleft\arraybackslash}p{(\linewidth - 6\tabcolsep) * \real{0.2667}}@{}}
\toprule\noalign{}
\begin{minipage}[b]{\linewidth}\raggedright
Measurement
\end{minipage} & \begin{minipage}[b]{\linewidth}\raggedleft
Observed σ (ms)
\end{minipage} & \begin{minipage}[b]{\linewidth}\raggedleft
n
\end{minipage} & \begin{minipage}[b]{\linewidth}\raggedleft
SE = σ/√n (ms)
\end{minipage} \\
\midrule\noalign{}
\endhead
\bottomrule\noalign{}
\endlastfoot
E15 \texttt{frame\_interval\_ms} (parallel, post-eDMA) & 1.8 & 20 &
0.40 \\
E16 \texttt{total\_frame\_ms} (alternating, pre-fix) & 32.5 & 40 (20 per
direction × 2) & 5.1 aggregate, 0.5 per direction \\
E16 \texttt{total\_frame\_ms} (alternating, post-Fix B) & 2.5 & 40 &
0.40 \\
E18 sweep A/B \texttt{total\_frame\_ms} & 2.2 & 40 & 0.35 \\
E18 sweep C \texttt{total\_frame\_ms} & 2.5 & 40 & 0.40 \\
\end{longtable}
}

Reported effect sizes are 9--83 ms. SE is \textasciitilde0.4 ms. The
ratio of effect to standard error is therefore 20--200 ×, so a larger
sample would not change any headline conclusion. The x20 choice for
E14/E15 and x40 for E16/E18 is a pragmatic trade-off: large enough that
the confidence interval is comfortably narrower than the reported
differences, small enough that a session completes in 3--10 s of wall
time (relevant because REPL-driven experiments are more likely to suffer
a transient fault --- CDC-ACM stall, LFS-busy race --- the longer they
run).

We did not run a formal power analysis a priori because the effect sizes
were unknown before the first measurements; the x20 choice was reactive
(20 is the default \texttt{repetitions} for the parallel pipeline in the
early diag modules) and the x40 choice was a conscious upgrade once the
alternating sweeps were introduced and we wanted ≥ 19 samples per
direction after the first-drop and the alternation split.

\begin{center}\rule{0.5\linewidth}{0.5pt}\end{center}

\section{2. Why the first sample is
dropped}\label{why-the-first-sample-is-dropped}

The first iteration inside a measurement loop differs from subsequent
iterations in three ways:

\begin{enumerate}
\def\labelenumi{\arabic{enumi}.}
\tightlist
\item
  \textbf{Switch-count parity}: E16/E18 alternating starts on one
  direction (cam\_a by convention) before any flip has happened in the
  measured loop. Its ``drain'' path is therefore different from the
  drain paths of iterations 1 \ldots{} N−1.
\item
  \textbf{DMA queue occupancy}: the warm-up frame leaves buffers queued;
  the first measurement iteration consumes those. Subsequent iterations
  see a full steady-state queue at entry.
\item
  \textbf{TPU package-cache miss}: the first \texttt{invoke()} after a
  model load, or after the pipeline has been stopped and restarted, pays
  a one-time package-cache walk that later invokes do not.
\end{enumerate}

Dropping the first sample removes all three confounds with a cost of n−1
instead of n.~At n = 40, that is a 2.5 \% data loss, worth trading for
unbiased steady-state statistics. Note that \emph{none} of the three
confounds above could be corrected by a longer run; only by dropping the
outlier.

We have observed the first-sample elevation empirically: in every
\texttt{s045\_e18\_post\_refactor/003\_e18\_C\_alt\_cam0\_cam1.csv} run,
the first row's \texttt{total\_frame\_ms} is \textasciitilde15 ms above
the cohort mean. The appendix generator drops it automatically
(\href{../experiments/_build_appendix.py}{\texttt{../experiments/\_build\_appendix.py}}
in \texttt{e16\_session} and \texttt{e18\_session}).

\begin{center}\rule{0.5\linewidth}{0.5pt}\end{center}

\section{3. Bessel-corrected σ, not population
σ}\label{bessel-corrected-ux3c3-not-population-ux3c3}

We report \texttt{statistics.stdev} (Bessel-corrected sample standard
deviation, denominator \texttt{n−1}) rather than
\texttt{statistics.pstdev} (population, denominator \texttt{n}). With n
= 39--40 the difference is \textless{} 2 \% and does not affect any
conclusion, but Bessel's correction is the unbiased estimator when the
sample is drawn from a larger notional population (e.g.~a bootstrapped
rerun), which is the situation we are in: each measurement run is one
sample from the distribution of ``runs on this hardware with these
parameters''.

\begin{center}\rule{0.5\linewidth}{0.5pt}\end{center}

\section{4. Mean vs median}\label{mean-vs-median}

The headline per-stage tables in \url{evaluation.md} use arithmetic
means because:

\begin{enumerate}
\def\labelenumi{\arabic{enumi}.}
\tightlist
\item
  The distributions are close to symmetric once the first sample is
  dropped (skewness below 0.5 on all measured stages post-warm-up).
\item
  Arithmetic means commute with the \texttt{1000/mean\_ms} FPS
  derivation --- reporting median and then deriving ``median FPS'' would
  be a different, more complicated statistic that did not match the
  reader's intuition of ``average throughput''.
\item
  Outliers in the archived CSVs are rare (\textless{} 1 \%) and when
  they appear they are either a single \texttt{max} above the cohort by
  \textasciitilde15 ms
  (e.g.~\texttt{s043/003\_e18\_C\_alt\_cam0\_cam1.csv} row 38 with total
  154 ms vs cohort mean 145 ms, consistent with a CDC-ACM transient) or
  a cluster at the end of a session where a sensor-side wake effect
  lagged. Either way the contribution to the mean is well below the
  cross-session noise band.
\end{enumerate}

Where we \emph{do} report median --- the per-direction split in
alternating sweeps (cam0 vs cam1) --- we do so because the data is
bimodal and the mean of the combined cohort is not meaningful. The
appendix tables in \url{evaluation.md} §4 split cam0 and cam1 into
separate unimodal sub-tables for exactly this reason.

\begin{center}\rule{0.5\linewidth}{0.5pt}\end{center}

\section{5. No p-values, no confidence intervals on derived
statistics}\label{no-p-values-no-confidence-intervals-on-derived-statistics}

We do not report p-values, t-tests, or ANOVA. Reasons:

\begin{enumerate}
\def\labelenumi{\arabic{enumi}.}
\tightlist
\item
  Effect sizes are all ≥ 10 × the standard error, so any standard test
  would return p \textless{} 10⁻¹⁰ --- not informative.
\item
  The hypotheses we are testing are directional and pre-registered in
  the form of the implementation chapters (e.g. ``the eDMA memcpy should
  save 9 ms'', ``Fix B should eliminate the directional asymmetry'');
  after-the-fact null-hypothesis testing adds nothing.
\item
  The \textbf{cross-session agreement} in Table 4 of \url{evaluation.md}
  is the rigorous claim: three independent measurements of the same
  quantity agree within ≤ 1 ms. That is stronger evidence than any
  single-session confidence interval, because it tests the stability of
  the measurement pipeline --- not just the within-session variance.
\end{enumerate}

Where a single number matters \emph{in isolation} (e.g.~the 83 ms per-
switch tax), we quote it alongside its σ and its n so a reader can
compute any confidence interval they care about.

\begin{center}\rule{0.5\linewidth}{0.5pt}\end{center}

\section{\texorpdfstring{6. FPS derivation --- why
\texttt{1000\ /\ mean(total\_ms)} instead of direct
FPS}{6. FPS derivation --- why 1000 / mean(total\_ms) instead of direct FPS}}\label{fps-derivation-why-1000-meantotal_ms-instead-of-direct-fps}

Every reported FPS is derived from the mean of per-iteration
\texttt{total\_frame\_ms}, not measured as ``frames in a fixed time
window''. The equivalence holds under the following condition:

\begin{quote}
If per-iteration times are i.i.d. with mean μ and finite variance, then
the long-run frames-per-second observed over a window of length T
converges almost surely to 1000/μ as T → ∞.
\end{quote}

Our iteration times satisfy i.i.d.-after-warm-up-drop within session.
The σ (1--3 ms on a 60--150 ms mean) is well within the regime where the
equivalence holds numerically: 1000/mean differs from a window- count
FPS by \textless{} 0.1 \%. See any standard text on renewal processes;
the 1000/mean derivation is the pointwise estimator.

We use the derivation because it lets us report FPS for a single
measurement window (a sweep), not just a long sustained run, and it
surfaces the per-stage contributions directly --- the user can read off
``the reason I am at 16 FPS is that
\texttt{invoke\ +\ to\_tensor\ +\ detect\ =\ 62\ ms}'' rather than
trying to attribute a windowed frame count.

\begin{center}\rule{0.5\linewidth}{0.5pt}\end{center}

\section{7. Handling of sessions with n \textless{}
20}\label{handling-of-sessions-with-n-20}

Sessions \href{../experiments/s016_e15_512/}{s016} through
\href{../experiments/s020_e15_512/}{s020} and
\href{../experiments/s027_e15_512/}{s027} through
\href{../experiments/s030_e15_512/}{s030} contain single-shot runs with
n = 1 each --- they were pre/post-eDMA ``does it still work'' sanity
probes during the bring-up phase and do not provide a defensible mean.
These sessions are \textbf{excluded from every statistical comparison}
in this paper. They appear in Appendix C of
\href{../experiments/README.md}{experiments/README.md} for completeness
but their ``FPS'' rows are derived from a single data point and
explicitly flagged as such.

Sessions \href{../experiments/s036_e17_drain_ab/}{s036},
\href{../experiments/s037_e17_drain_ab/}{s037}, and
\href{../experiments/s040_e17_drain_ab/}{s040} are empty sessions where
the REPL driver failed mid-run and no experiment completed. They are
present in the archive so the session-ID sequence stays contiguous; they
contribute nothing to any result in this paper.

\begin{center}\rule{0.5\linewidth}{0.5pt}\end{center}

\section{8. What the reader should
check}\label{what-the-reader-should-check}

A reviewer who wants to stress-test the statistical choices in this
paper can re-run the appendix generator against the archived CSVs and
confirm that:

\begin{enumerate}
\def\labelenumi{\arabic{enumi}.}
\tightlist
\item
  Every headline mean in the paper tables matches the regenerated
  appendix to within rounding.
\item
  The σ values reported are Bessel-corrected (they match the
  \texttt{statistics.stdev} function applied to the raw columns with the
  first sample dropped).
\item
  The cross-session agreement in Table 4 of \url{evaluation.md} holds
  with any reasonable choice of warm-up-drop length (0, 1, or 2 samples)
  --- the claim is robust, not cherry-picked.
\end{enumerate}

The command to run is in \url{artifact.md} §7 ``Data integrity check''.
